% ======================================================================
\section{The observable and its expansion}\label{sec:observable}
% ======================================================================

\subsection{Single inclusive gluon spectrum}\label{sec:observable_def}

The projectile, before the collision, is described by its valence state $\ket{\Psi_{\rm v}}$ times the dressing of the soft modes, $\Omega\ket{0}$, where $\ket0$ is the vacuum of the soft modes. The target is described by the eikonal scattering matrix $S$. At high energy the interaction of a soft gluon and of the valence charges with the target is a colour rotation by the adjoint Wilson line $U(\v x)$ at fixed transverse position:
\begin{align}\label{eq:Srotation}
S^\dagger\rho^a(\v x)S&=U^{ab}(\v x)\,\rho^b(\v x)\,,\notag\\
S^\dagger a^a_i(k^+,\v x)S&=U^{ab}(\v x)\,a^b_i(k^+,\v x)\,,\qquad
S^\dagger a^{\dagger a}_i(k^+,\v x)S=U^{ab}(\v x)\,a^{\dagger b}_i(k^+,\v x)\,.
\end{align}
After the collision the state is $S\,\Omega\ket0$. A gluon observed in the final state is a dressed (physical) gluon; the number of dressed gluons is counted by the bare number operator in the basis rotated back by $\Omega^\dagger$~\cite{Kovner:2006wr}. The single inclusive spectrum is therefore
\begin{align}\label{eq:observable}
\frac{d^3N}{d^3k}=\frac{1}{(2\pi)^3}\int_{w,w'}e^{-i\v k\cdot(\v w-\v w')}\,
\Big\langle\bra0\,\Omega^\dagger S^\dagger\Omega\;a^{\dagger a}_i(k^+,\v w')\,a^a_i(k^+,\v w)\;\Omega^\dagger S\,\Omega\ket0\Big\rangle\,,
\end{align}
where $\langle\dots\rangle$ denotes the average over the valence state, i.e.\ over the projectile charge configurations, taken at the end. The factor $(2\pi)^{-3}$ and the phase follow from $a^{\dagger a}_i(k)a^a_i(k)=\int_{w,w'}e^{-i\v k\cdot(\v w-\v w')}a^{\dagger a}_i(k^+,\v w')a^a_i(k^+,\v w)$ and the normalisation~\eqref{eq:aadagger}. The amplitude stands at $\v w$ and the conjugate amplitude at $\v w'$; we use this convention, $e^{-i\v k\cdot(\v w-\v w')}$ with the conjugated factor at the primed points, throughout the paper.

Inserting $SS^\dagger=1$ and defining $\bar\Omega\equiv S^\dagger\Omega S$, which is $\Omega$ with $\rho\to U\rho$ and $a\to Ua$ according to Eq.~\eqref{eq:Srotation}, we obtain
\begin{align}
\frac{d^3N}{d^3k}=\frac{1}{(2\pi)^3}\int_{w,w'}e^{-i\v k\cdot(\v w-\v w')}\,\Big\langle\bra0\,\Omega^\dagger\bar\Omega\;U^{ab'}(\v w')a^{\dagger b'}_i(k^+,\v w')\,U^{ab}(\v w)a^b_i(k^+,\v w)\;\bar\Omega^\dagger\Omega\ket0\Big\rangle\,.
\end{align}
A bar on any coefficient of $\Omega$ denotes, here and below, the replacement $\rho^a(\v u)\to U^{ab}(\v u)\rho^b(\v u)$ of every charge in it. Between $a(k^+,\v w)$ and $\bar\Omega^\dagger\Omega\ket0$ we insert a complete set of soft states. The vacuum does not contribute, and states with three or more gluons contribute at order $g^6$. Hence
\begin{align}\label{eq:N_amplitudes}
\frac{d^3N}{d^3k}=\frac{1}{(2\pi)^3}\Big\langle\mathcal A^{a\dagger}_i(k)\,\mathcal A^a_i(k)+\int dq^+\!\int_{v}\mathcal A^{ab\dagger}_{ij}(k;q^+,\v v)\,\mathcal A^{ab}_{ij}(k;q^+,\v v)\Big\rangle\,,
\end{align}
with the one-gluon and two-gluon amplitudes
\begin{align}\label{eq:amplitudes_def}
\mathcal A^a_i(k)&=\sqrt{2\pi}\int_{w}e^{-i\v k\cdot\v w}\,U^{ab}(\v w)\,\bra{g^b_i(k^+,\v w)}\bar\Omega^\dagger\Omega\ket0\,,\notag\\
\mathcal A^{ab}_{ij}(k;q^+,\v v)&=\sqrt{2\pi}\int_{w}e^{-i\v k\cdot\v w}\,U^{ac}(\v w)\,\bra{g^b_j(q^+,\v v)\,g^c_i(k^+,\v w)}\bar\Omega^\dagger\Omega\ket0\,.
\end{align}
The colour index $b$ of the second gluon in $\mathcal A^{ab}_{ij}$ is that of the bare state; since the second gluon is integrated over, its rotation by $U(\v v)$ drops out of the sum. The operators $\bar\Omega^\dagger\Omega$ are evaluated with the expansion~\eqref{eq:OmegaX}: $\Omega$ acting on $\ket0$ involves only $\mf N$, $\mf A_2$ and $\mf B_2$; $\bar\Omega^\dagger$ then acts on states with up to two gluons. Note that $\bar\Omega^\dagger\Omega=1$ for $U=1$, so that every amplitude vanishes in the absence of the target.

From here on $\mf A\equiv\mf A^{(1)}_2$ and $\mf A^\dagger\equiv\mf A^{(1)\dagger}_2$; the unitarity relation $\mf A^{(1)\dagger}_1(p^+,\v x)=-\mf A(p^+,-\v x)$ is used to eliminate $\mf A^{(1)}_1$ wherever possible, and $\mf A^{(3)}\equiv\mf A^{(3)}_2$.

\subsection{Leading order}\label{sec:LO}

At order $g$ only $\mf A$ contributes:
\begin{align}\label{eq:ALO}
\mathcal A^{a}_{{\rm LO}\,i}(k)&=\int_{w}e^{-i\v k\cdot\v w}\Big[U^{ab}(\v w)\,\mf A^b_i(k^+,-\v w)-\bar{\mf A}^a_i(k^+,-\v w)\Big]
\notag\\&=\frac{ig}{\sqrt2\,\pi\sqrt{k^+}}\int_{w,x}e^{-i\v k\cdot\v w}\,\Kv^i(\v x-\v w)\big[U(\v w)-U(\v x)\big]^{ab}\rho^b(\v x)\,.
\end{align}
The first term is the emission of the gluon before the collision, followed by the scattering of the gluon; the second is the emission after the collision by the rotated charge. The LO spectrum is
\begin{align}\label{eq:LO}
\frac{d^3N_{\rm LO}}{d^3k}=\ &\frac{1}{(2\pi)^3}\,\frac{g^2}{2\pi^2}\,\frac{1}{k^+}\int_{x,x',w,w'}e^{-i\v k\cdot(\v w-\v w')}\,
\Kv(\v x-\v w)\cdot\Kv(\v x'-\v w')
\notag\\&\times
\Big\langle\big[U(\v x')-U(\v w')\big]^{ab'}\rho^{b'}(\v x')\,\big[U(\v x)-U(\v w)\big]^{ab}\rho^{b}(\v x)\Big\rangle\,.
\end{align}
Note that $\frac{1}{(2\pi)^3}\frac{g^2}{2\pi^2}=\frac{\alpha_s}{4\pi^4}$.

\paragraph{Colour-singlet projectile.} For a projectile that is a colour singlet, or more generally whose charge correlator is colour diagonal,
\begin{align}\label{eq:singlet}
\big\langle\rho^a(\v x)\rho^b(\v y)\big\rangle=\frac{\delta^{ab}}{N_c^2-1}\,n(\v x,\v y)\,,\qquad n(\v x,\v y)=\big\langle\rho^f(\v x)\rho^f(\v y)\big\rangle\,,
\end{align}
and the LO spectrum becomes
\begin{align}\label{eq:LOsinglet}
\frac{d^3N_{\rm LO}}{d^3k}=\ &\frac{\alpha_s}{4\pi^4(N_c^2-1)}\,\frac{1}{k^+}\int_{x,y,w,w'}e^{-i\v k\cdot(\v w-\v w')}\,\Kv(\v x-\v w)\cdot\Kv(\v y-\v w')
\notag\\&\times\,n(\v x,\v y)\,
\mathcal O_{\rm LO}(\v x,\v y;\v w,\v w')\,,
\end{align}
with the target operator
\begin{align}\label{eq:OLO}
\mathcal O_{\rm LO}(\v x,\v y;\v w,\v w')={\rm Tr}\Big\{\big[U^\dagger(\v x)-U^\dagger(\v w)\big]\big[U(\v y)-U(\v w')\big]\Big\}\,.
\end{align}
Since $U$ is real, ${\rm Tr}\{A^\dagger B\}={\rm Tr}\{B^\dagger A\}$ for adjoint matrices, and for a symmetric $n$ the integrand is invariant under $(\v x,\v w)\leftrightarrow(\v y,\v w')$; in the colour-singlet form the sign of the exponent is therefore immaterial. Eq.~\eqref{eq:LOsinglet} is the formula of Kovchegov and Tuchin~\cite{Kovchegov:2001sc}: for a single quark, $n\to C_F$ times a $\delta$-function, it reproduces their normalisation $\alpha_sC_F/(4\pi^4)$ exactly. It is also the operator that appears in the LL analysis of Kovner and Lublinsky~\cite{Kovner:2006wr,Kovner:2006ge}.

\paragraph{Momentum space and $k_T$ factorisation.} With $f(\v k)=\int_xe^{i\v k\cdot\v x}f(\v x)$ and Eq.~\eqref{eq:FT_K} of Appendix~\ref{app:fourier}, the LO amplitude takes the form
\begin{align}\label{eq:ALOmom}
\mathcal A^{a}_{{\rm LO}\,i}(k)=\frac{\sqrt2\,g}{\sqrt{k^+}}\int\frac{d^2q}{(2\pi)^2}\,\Big[\frac{k^i}{\v k^2}-\frac{q^i}{\v q^2}\Big]\,U^{ab}(\v q-\v k)\,\rho^b(-\v q)\,,\qquad
U^{ab}(\v l)=\int_xe^{i\v l\cdot\v x}U^{ab}(\v x)\,,
\end{align}
so that the gluon is produced with the Lipatov-type vertex $\v k/\v k^2-\v q/\v q^2$, where $\v q$ is the momentum supplied by the projectile and $\v k-\v q$ the momentum transferred by the target. The term $U^{ab}(\v l)\propto\delta^{ab}\delta^{(2)}(\v l)$ of a trivial target does not contribute, since the vertex vanishes at $\v q=\v k$. For a dilute target, $U^{ab}(\v x)=\delta^{ab}+f^{acb}\alpha^c(\v x)+O(\alpha^2)$, the amplitude is linear in $\alpha$, and with $\big|\frac{\v k}{\v k^2}-\frac{\v q}{\v q^2}\big|^2=\frac{(\v k-\v q)^2}{\v k^2\v q^2}$ the spectrum reduces to the $k_T$-factorised form~\cite{Gribov:1983ivg,Kovchegov:2001sc,Kharzeev:2003wz,Blaizot:2004wu}. For a projectile and a target that are homogeneous in the transverse plane, with $\langle\rho^{b'}(\v q')\rho^b(-\v q)\rangle=\delta^{bb'}(2\pi)^2\delta^{(2)}(\v q-\v q')P(\v q)$ and $\langle\alpha^{c}(\v l)\alpha^{d}(-\v l')\rangle=\delta^{cd}(2\pi)^2\delta^{(2)}(\v l-\v l')T(\v l)$,
\begin{align}\label{eq:kTfact}
\frac{d^3N_{\rm LO}}{d^3k}\simeq\frac{\alpha_s N_c(N_c^2-1)}{\pi^2}\,\frac{S_\perp}{k^+}\int\frac{d^2q}{(2\pi)^2}\,\frac{(\v k-\v q)^2}{\v k^2\,\v q^2}\,P(\v q)\,T(\v k-\v q)\,,
\end{align}
where $S_\perp$ is the transverse area.

\subsection{Next-to-leading order: amplitudes}\label{sec:NLOamplitudes}

At order $g^3$ the one-gluon amplitude receives five types of contributions. Collecting the terms of $\bar\Omega^\dagger\Omega$ that connect $\ket0$ with $\bra{g}$ at this order we find
\begin{align}\label{eq:A3amp}
\mathcal A^a_{(3)\,i}(k)=\int_w e^{-i\v k\cdot\v w}\Big\{\mathcal A_{N}+\mathcal A_{A^{(3)}}+\mathcal A_{B_3}+\mathcal A_{AB_2}+\mathcal A_{CB_2}\Big\}^a_i\,,
\end{align}
with
\begin{align}
\mathcal A_{N}&=-\bar{\mf A}^a_i(k^+,-\v w)\,\mf N_2+U^{ab}(\v w)\,\bar{\mf N}_2\,\mf A^b_i(k^+,-\v w)-\big[\bar{\mf N}_2,\bar{\mf A}^a_i(k^+,-\v w)\big]\,,
\notag\\
\mathcal A_{A^{(3)}}&=U^{ab}(\v w)\,\mf A^{(3)b}_i(k^+,-\v w)-\bar{\mf A}^{(3)a}_i(k^+,-\v w)\,,
\notag\\
\mathcal A_{B_3}&=\int\frac{dp^+}{2\pi}\int_x\bar{\mf B}^{\dagger ba}_{3ji}(k^+,\v w;p^+,-\v x)\Big[U^{bc}(\v x)\,\mf A^c_j(p^+,-\v x)-\bar{\mf A}^b_j(p^+,-\v x)\Big]\,,
\notag\\
\mathcal A_{AB_2}&=2\int\frac{dp^+}{2\pi}\int_x\Big[U^{ad}(\v w)U^{bc}(\v x)\,\bar{\mf A}^{\dagger b}_j(p^+,-\v x)\,\mf B^{cd}_{2ji}(k^+,-\v w;p^+,-\v x)
\notag\\&\hspace{2.2cm}
-\bar{\mf A}^{\dagger b}_j(p^+,-\v x)\,\bar{\mf B}^{ba}_{2ji}(k^+,-\v w;p^+,-\v x)\Big]\,,
\notag\\
\mathcal A_{CB_2}&=2\int\frac{dp^+}{2\pi}\int_{x,z}\Big[\bar{\mf C}^{\dagger dca}_{1jki}(k^+,\v w-\v z;p^+,\v z-\v x)\,U^{db}(\v x)U^{ce}(\v z)\,\mf B^{be}_{2jk}(k^+-p^+,-\v z;p^+,-\v x)
\notag\\&\hspace{2.2cm}
-\bar{\mf C}^{\dagger bca}_{1jki}(k^+,\v w-\v z;p^+,\v z-\v x)\,\bar{\mf B}^{bc}_{2jk}(k^+-p^+,-\v z;p^+,-\v x)\Big]\,.
\end{align}
The term $\mathcal A_N$ contains the normalisation: the term $\bar{\mf A}^{(3)}_1$ of $\bar\Omega^\dagger$ contributes, through the unitarity relation~\eqref{eq:A31unit}, the commutator $-[\bar{\mf N}_2,\bar{\mf A}]$. $\mathcal A_{A^{(3)}}$ is the $O(g^3)$ one-gluon amplitude of the wave function, before and after the collision. $\mathcal A_{B_3}$ describes the scattering of an already emitted gluon of momentum $p^+$ into the observed gluon by its interaction with the valence charges after the collision ($\bar{\mf B}_3^\dagger$); $\mathcal A_{AB_2}$, the absorption, after the collision, of one of the two gluons of the two-gluon component of the wave function; $\mathcal A_{CB_2}$, the recombination of the two gluons of that component into the observed one. In each term the two contributions correspond to an interaction of the soft gluons with the target before and after the respective valence interaction; their difference vanishes for $U=1$.

The terms of $\mathcal A_N$ are not all of the same order in $\rho$. With the ordering of the three-$\rho$ part of $\mf A^{(3)}$ derived in Section~\ref{sec:coefficients}, $[\mf A^{(3)}]_{\rho\rho\rho}=\mf A\,\mf N_2$, the commutator in $\mathcal A_N$ cancels against the corresponding part of $\mathcal A_{A^{(3)}}$, and the four-$\rho$ (normalisation) terms of the amplitude add up to
\begin{align}\label{eq:Nterms}
-\bar{\mf A}\,\mf N_2+U\,\bar{\mf N}_2\,\mf A+U\,\mf A\,\mf N_2-\bar{\mf N}_2\,\bar{\mf A}\,,
\end{align}
which vanishes for $U=1$, as it must. With the opposite ordering, $\mf N_2\mf A$, a spurious two-$\rho$ term $\propto\log(\vee/\Lambda)$ would remain.

The two-gluon amplitude at order $g^2$ is
\begin{align}\label{eq:A2gamp}
\mathcal A^{ab}_{ij}(k;q^+,\v v)=\frac{2}{\sqrt{2\pi}}\int_w e^{-i\v k\cdot\v w}\Big\{\mathcal A_B+\mathcal A^{(1)}_{AA}+\mathcal A^{(2)}_{AA}-\mathcal A_{CA}\Big\}^{ab}_{ij}\,,
\end{align}
with
\begin{align}
\mathcal A_B&=U^{ac}(\v w)\,\mf B^{bc}_{2ji}(k^+,-\v w;q^+,-\v v)-U^{\dagger bc}(\v v)\,\bar{\mf B}^{ca}_{2ji}(k^+,-\v w;q^+,-\v v)\,,
\notag\\
\mathcal A^{(1)}_{AA}&=\tfrac12\Big[U^{\dagger bc}(\v v)\,\bar{\mf A}^c_j(q^+,-\v v)\,\bar{\mf A}^a_i(k^+,-\v w)-U^{ac}(\v w)\,\bar{\mf A}^d_j(q^+,-\v v)\,U^{\dagger bd}(\v v)\,\mf A^c_i(k^+,-\v w)\Big]\,,
\notag\\
\mathcal A^{(2)}_{AA}&=\tfrac12\Big[U^{\dagger bc}(\v v)\,\bar{\mf A}^a_i(k^+,-\v w)\,\bar{\mf A}^c_j(q^+,-\v v)-\bar{\mf A}^a_i(k^+,-\v w)\,\mf A^b_j(q^+,-\v v)\Big]\,,
\notag\\
\mathcal A_{CA}&=\int_z\Big[\bar{\mf C}^{dac}_{1jik}(q^++k^+,\v z-\v w;q^+,\v w-\v v)\,U^{\dagger bd}(\v v)\,U^{ce}(\v z)\,\mf A^e_k(q^++k^+,-\v z)
\notag\\&\hspace{1.2cm}
-U^{\dagger bd}(\v v)\,\bar{\mf C}^{dac}_{1jik}(q^++k^+,\v z-\v w;q^+,\v w-\v v)\,\bar{\mf A}^c_k(q^++k^+,-\v z)\Big]\,.
\end{align}
$\mathcal A_B$ is the two-gluon component of the wave function scattered on the target; $\mathcal A^{(1)}_{AA}$ and $\mathcal A^{(2)}_{AA}$ arise from the product of two single emissions, one before and one after the collision, in the two operator orderings generated by $\bar\Omega^\dagger$; and $\mathcal A_{CA}$ is the splitting, after the collision, of a gluon of momentum $q^++k^+$ emitted before the collision into the observed gluon and a second gluon of momentum $q^+$.

\subsection{Next-to-leading order: Groups and rows}\label{sec:groups}

The $O(g^4)$ cross section is
\begin{align}\label{eq:NLOdecomp}
\frac{d^3N}{d^3k}=\frac{d^3N_{\rm LO}}{d^3k}+\frac{d^3N^{\rm(I)}}{d^3k}+\frac{d^3N^{\rm(II)}}{d^3k}+\frac{d^3N^{\rm(III)}}{d^3k}+O(g^6)\,,
\end{align}
where
\begin{align}\label{eq:Groups}
\frac{d^3N^{\rm(I)}}{d^3k}&=\frac{1}{(2\pi)^3}\Big\langle\mathcal A^{a\dagger}_{{\rm LO}\,i}(k)\,\mathcal A^a_{(3)\,i}(k)\Big\rangle+{\rm c.c.}\,,\notag\\
\frac{d^3N^{\rm(II)}}{d^3k}&=\frac{1}{(2\pi)^3}\int dq^+\!\int_v\Big\langle\big|\mathcal A_B\big|^2+\big|\mathcal A^{(1)}_{AA}\big|^2+\big|\mathcal A^{(2)}_{AA}\big|^2+\big|\mathcal A_{CA}\big|^2\Big\rangle\,,\notag\\
\frac{d^3N^{\rm(III)}}{d^3k}&=\frac{1}{(2\pi)^3}\int dq^+\!\int_v\Big\langle\mathcal A^{(1)\dagger}_{AA}\mathcal A_B+\mathcal A^{(2)\dagger}_{AA}\mathcal A_B-\mathcal A^{\dagger}_{CA}\mathcal A_B
\notag\\&\hspace{3.3cm}+\mathcal A^{(2)\dagger}_{AA}\mathcal A^{(1)}_{AA}-\mathcal A^{\dagger}_{CA}\mathcal A^{(1)}_{AA}-\mathcal A^{\dagger}_{CA}\mathcal A^{(2)}_{AA}\Big\rangle+{\rm c.c.}
\end{align}
Here $|\mathcal A|^2\equiv\mathcal A^{ab\dagger}_{ij}\mathcal A^{ab}_{ij}$, the sign of $\mathcal A_{CA}$ is that of Eq.~\eqref{eq:A2gamp}, and the momentum of the integrated gluon is renamed $q^+\to p^+$ below. Group~I is the interference of the LO amplitude with the $O(g^3)$ one-gluon amplitude; Group~II consists of the squares of the four components of the two-gluon amplitude; Group~III of their interferences. Groups~I and~III are written with their complex conjugates; Group~II is Hermitian by itself. We call the individual products of coefficients \emph{rows}. They are listed in Table~\ref{tab:rows}, together with the range of the longitudinal integration, the number of charges and the parts of the final result to which they contribute.

\begin{table}[t]
\centering
\renewcommand{\arraystretch}{1.2}
\footnotesize
\begin{tabular}{@{}l p{4.0cm} p{2.9cm} l p{3.4cm}@{}}
\toprule
Row & product (Eqs.~\eqref{eq:A3amp},~\eqref{eq:A2gamp}) & range of $p^+$ & $\#\rho$ & contributes to\\
\midrule
I.I & LO$^\dagger\times$ normalisation terms~\eqref{eq:Nterms} & --- & 4 & 4$\rho$ ($\log\frac{\vee}{\Lambda}$)\\
I.II & LO$^\dagger\times[U\mf A^{(3)}-\bar{\mf A}^{(3)}]$ & --- & 2--4 & 4$\rho$, BFKL, JIMWLK, 3$\rho$, finite\\
I.III & LO$^\dagger\times\bar{\mf B}_3^\dagger\mf A$ & $\Lambda<p^+<\vee$, $|p^+-k^+|>\Lambda$ & 3, 4 & BFKL, 3$\rho$, finite\\
I.IV & LO$^\dagger\times\bar{\mf A}^\dagger\mf B_2$ & $\Lambda<p^+<\vee-k^+$ $[\vee]$ & 3, 4 & 4$\rho$, BFKL, 3$\rho$; finite part $=-$III.I\\
I.V & LO$^\dagger\times\bar{\mf C}_1^\dagger\mf B_2$ & $\Lambda<p^+<k^+-\Lambda$ & 2, 3 & JIMWLK, 3$\rho$, finite\\
\midrule
II.I & $|\mathcal A_B|^2$ & $\Lambda<p^+<\vee-k^+$ $[\vee]$ & 2--4 & 4$\rho$, BFKL, JIMWLK, 3$\rho$, finite\\
II.II & $|\mathcal A^{(1)}_{AA}|^2$ & $\Lambda<p^+<\vee$ & 4 & 4$\rho$\\
II.III & $|\mathcal A^{(2)}_{AA}|^2$ & $\Lambda<p^+<\vee$ & 4 & 4$\rho$\\
II.IV & $|\mathcal A_{CA}|^2$ & $\Lambda<p^+<\vee-k^+$ & 2 & BFKL, JIMWLK, DGLAP, RC, finite\\
\midrule
III.I & $\mathcal A^{(1)\dagger}_{AA}\mathcal A_B$ & $\Lambda<p^+<\vee-k^+$ $[\vee]$ & 3, 4 & 4$\rho$, BFKL, 3$\rho$; finite part $=-$I.IV\\
III.II & $\mathcal A^{(2)\dagger}_{AA}\mathcal A_B$ & $\Lambda<p^+<\vee-k^+$ $[\vee]$ & 3, 4 & 4$\rho$, BFKL, 3$\rho$, finite\\
III.III & $-\mathcal A^{\dagger}_{CA}\mathcal A_B$ & $\Lambda<p^+<\vee-k^+$ & 2, 3 & BFKL, JIMWLK, 3$\rho$, finite\\
III.IV & $\mathcal A^{(2)\dagger}_{AA}\mathcal A^{(1)}_{AA}$ & $\Lambda<p^+<\vee$ & 4 & 4$\rho$\\
III.V & $-\mathcal A^{\dagger}_{CA}\mathcal A^{(1)}_{AA}$ & $\Lambda<p^+<\vee-k^+$ & 3 & BFKL, 3$\rho$, finite\\
III.VI & $-\mathcal A^{\dagger}_{CA}\mathcal A^{(2)}_{AA}$ & $\Lambda<p^+<\vee-k^+$ & 3 & BFKL, 3$\rho$, finite\\
\bottomrule
\end{tabular}
\caption{The rows of the $O(g^4)$ cross section. ``LO$^\dagger$'' is the conjugate LO amplitude, Eq.~\eqref{eq:ALO}; $p^+$ is the longitudinal momentum of the second (integrated or virtual) gluon. Where two ranges apply, the first refers to the part linear in $\rho$ of $\mf B_2$ (its parent gluon carries $p^++k^+$) and the bracket $[\vee]$ to its $\{\rho,\rho\}$ part, $\Lambda<p^+<\vee$. ``$\#\rho$'' is the number of charges before reordering; ``4$\rho$'' and ``3$\rho$'' denote the terms quartic and cubic in $\rho$ multiplying a rapidity logarithm, which cancel (Section~\ref{sec:rapidity}). Row I.II contains in addition the ultraviolet $\beta_0$ terms of the vertex correction, which cancel against Row I.V (Section~\ref{sec:dglap_rc}).}
\label{tab:rows}
\end{table}

The longitudinal ranges in Table~\ref{tab:rows} follow from the kinematics of the coefficients. All soft gluons satisfy $\Lambda<p^+<\vee$. Whenever a soft gluon of momentum $p^++k^+$ appears as an intermediate state (the parent of $\mf C_1$, the part of $\mf B_2$ linear in $\rho$), the condition $p^++k^+<\vee$ applies. The $\{\rho,\rho\}$ part of $\mf B_2$ describes two independent emissions and has support on the full square $\Lambda<k^+,p^+<\vee$; this is required for the cancellation of the four-$\rho$ terms (Section~\ref{sec:fourrho}). In $\mf B_3$ the momentum difference of the incoming and outgoing gluon is carried by the valence charges and satisfies $|p^+-k^+|>\Lambda$.

The evaluation of each row consists of the colour algebra of the products of coefficients, with the charges kept in their order, the longitudinal integration over $p^+$, and the transverse integrations that can be done in closed form. Three kinds of terms emerge:
\begin{enumerate}
\item terms proportional to $\log(\vee/k^+)$ or $\log(k^+/\Lambda)$, quadratic, cubic or quartic in $\rho$, which arise from the regions $p^+\gg k^+$ and $p^+\ll k^+$ of the longitudinal integrals. They are discussed in Section~\ref{sec:rapidity};
\item terms with a collinear singularity, from the region in which the integrated gluon of Row~II.IV is collinear with the observed one, and terms proportional to $\beta_0$ (Section~\ref{sec:dglap_rc});
\item the remaining finite terms (Section~\ref{sec:finite}).
\end{enumerate}
Two rows require a remark. The non-logarithmic parts of Rows~I.IV and III.I cancel identically. Rows~II.II, II.III and III.IV are products of four coefficients $\mf A$, each linear in $\rho$; they contain only four-$\rho$ terms.
